\part{What are Floating Point Numbers?}
\frame{\partpage}
\begin{frame}[fragile]{How to create a floating point number?}
\begin{itemize}
\item by typing it, e.g. \pyth!a = 3.012! , or \pyth!a = 30.012e-1!, or \pyth!a=0.3012e+1! and other variants
\item by typecasting from integer
\begin{python}
n=1
a=float(n) # returns 1.0
\end{python}
\item as a result from a computation
\begin{python}
n=2
m=2
a=n/m   # returns 1.0
c=3+5j
\end{python}
\end{itemize}

We note \pyth!3.012 == 30.012e-1!. The decimal point "floats" depending on the given exponent.
\end{frame}
\begin{frame}[fragile]{Comparing floating point numbers}

Operations on floating point numbers rarely return the exact result:
\begin{python}
0.4 - 0.3 # returns 0.10000000000000003
\end{python}
This facts matters, when \emph{comparing floating point numbers}:
\begin{python}
0.4 - 0.3 == 0.1  # returns False
\end{python}

Facit:\\
Never compare floats by \pyth!==!. Ask instead if they are "near":
\begin{python}
abs((0.4-0.3) - 0.1 ) < 1.e-12
\end{python}
\end{frame}
\begin{frame}[fragile]{Internal representation}
Floats are represented by three quantities: the sign, the mantissa, and the exponent:
\[
\mathrm{sign} (x) \Bigl ( x_0 + x_1 \beta^{-1} + \ldots +
x_{t-1}\beta^{-(t-1)} \Bigr ) \beta^{e}
\]
with $\beta \in \mathbb{N}$ and
$x_0 \ne 0$, $0 \le x_i < \beta$.
\vfill
$x_0...x_{t-1}$ is the mantissa,\\
$\beta$ the basis,
$e$ exponent with $1-e_{\max} \le e \le e_{\max} $,
$t$ is the mantissa length.
\vfill
Normalization $x_0 \ne 0$  makes the representation unique\\
\hfill and saves in the case $\beta=2$ a bit.
\end{frame}
\begin{frame}[fragile]{Internal representation (Cont.)}
$0$ cannot be represented as a \it{normalized} float.
\vfill
Two zeros $+0$ and $-0$.
\vfill
The exponent has a special inicator to indicate when a number is not normalized.
\vfill
On your computer:
$\beta=2$, $t=52, e_{\max}=1023$. Together with the sign and the full range of the exponent this needs 64 bits.
\end{frame}
\begin{frame}[fragile]{Internal representation (Cont.)}
Smallest positive representable number is
\[
    \mathrm{fl}_{\min}=1.\;2^{-1023} \approx 10^{-308}
\]
and the largest is
\[
   \mathrm{fl}_{\max}=1.111\ldots1\;2^{1023} \approx 10^{308}.
\]
\vfill
Floating point numbers are not equally spaced in $[0,\mathrm{fl}_{max}]$.
\end{frame}
\begin{frame}[fragile]{Internal representation (Cont.)}
Gap at zero:\\
Distance between $0$ and the first positive number is $2^{-1023}$\\
Distance between the first and the second is smaller by a factor
$2^{-52} \approx 2.2\times 10^{-16}$,
\begin{center}
    \includegraphics[width=0.7\textwidth]{floatz}
\end{center}
\end{frame}
\begin{frame}[fragile]{Infinite and Not a Number}


There are in total
$
2 (\beta-1)\beta^{t-1}(2 e_{\max}+1)+1
$
floating point numbers.

Numbers bigger than $\mathrm{fl}_{\max}$ generate in numpy an \pyth$inf$
or an overflow exception

Numbers smaller than $\mathrm{fl}_{\max}$ generate subnormalized numbers.


\begin{python}
exp(1000.)  #  inf
a = inf
3 - a       # -inf
3 + a       #  inf

a + a  # inf
a - a  # nan
a / a  # nan      # not a number
\end{python}
\end{frame}
\begin{frame}[fragile]{Infinite and Not a Number (Cont.)}
Unexpected results from \pyth!nan!

\begin{python}
x = nan
x < 0 # False
x > 0 # False
x == x # False
\end{python}

The float \pyth{inf} behaves much more as expected:
\begin{python}
0 < inf # True
inf <= inf # True
inf == inf # True
-inf < inf # True
inf - inf  # nan
exp(-inf)  # 0
exp(1 / inf) # 1
\end{python}
\end{frame}
\begin{frame}[fragile]{Machine Epsilon}
\vfill
\begin{Definition}[machine epsilon]
The \emph{machine epsilon} or rounding unit
is the largest number $\varepsilon$ such that
\[
\mathrm{float}(1.0+\varepsilon)=1.0
\]
\end{Definition}
\vfill
Write a program which finds this number (or a good estimate for it)
\vfill
See also

\begin{python}
import sys
sys.float_info.epsilon # 2.220446049250313e-16 (depending on your system)
\end{python}

\end{frame}
